\documentclass[11pt]{article}
\usepackage[margin=1in]{geometry}
\usepackage{amsmath,amssymb}
\usepackage[hidelinks]{hyperref}
\providecommand{\citep}[1]{\cite{#1}}
\providecommand{\citet}[1]{\cite{#1}}
\usepackage{setspace}\onehalfspacing
\title{CEO Question-and-Answer Uncertainty and the Anticipation of Cash Acquisitions}
\author{}\date{}
\begin{document}
\maketitle
\section*{Abstract}

Firms host quarterly earnings calls and field analysts' unscripted questions, and the language a CEO uses in answering is increasingly read as a window into what managers know but have not yet disclosed. Yet when a firm is privately committed to an acquisition it has not announced, the executive holds the routine call unable to confirm or deny the one development that matters. Whether the spoken record tracks that withholding has not been characterized. We study the earnings-call transcripts and acquisition records of United States public firms from 2002 to 2018, covering 88,205 calls from 1,884 firms. From this language we construct a residual measure of CEO question-and-answer uncertainty---the part that remains once each executive's persistent speaking style is removed---and we track it around firms' cash acquisitions, with stock acquisitions as a comparison. We find that this residual uncertainty is elevated in the quarter before a cash acquisition, relative to the firm's own other quarters, with no comparable rise before stock acquisitions. Moreover, it becomes indistinguishable from zero once the deal is announced, even before completion, while the cash that funds the purchase stays on the balance sheet until the deal closes. Furthermore, the elevation concentrates in cash acquisitions rather than stock acquisitions, a difference that survives a formal pooled test. This pattern is not accounted for by the analyst cash-scrutiny we can measure; we read it as call language that tracks the deal's passage from private to public, a within-firm correlational regularity rather than a tested mechanism. Finally, the residual is also unrelated to the firm's post-call bid-ask spread, a standard gauge of information asymmetry, while the scripted presentation is, contemporaneously and on one-tailed within-firm tests, positively associated with it, consistent with the two parts of the call relating to the outside information environment in different ways. Overall, a constrained executive's unscripted answers appear to carry a readable, anticipatory trace of an acquisition the firm has not yet disclosed.

\section{1}

Earnings calls, and especially the unscripted question-and-answer session, are a setting in which managers' language carries information beyond the scripted presentation \citep{matsumoto2011}. Disclosure theory holds that an informed manager may rationally withhold private information when disclosure is costly \citep{verrecchia1983}, and that non-disclosure can persist in equilibrium because outside investors cannot distinguish an uninformed manager from one who is informed but silent \citep{dye1985}, so a withholding state is itself informative. A firm privately committed to an acquisition it has not announced sits in a concrete such state: information about preliminary merger negotiations can be material well before any definitive agreement \citep{basic_v_levinson}, and although the firm may stay silent it may not mislead once it speaks \citep{rule_10b5, basic_v_levinson}, so the executive can neither confirm nor deny the pending deal and is left to answer around it. When a manager cannot answer directly, the constraint can surface in the words themselves, for non-answers are common and not informationally empty: in a large share of calls managers leave requested information undisclosed, and, in the words of that literature, silence speaks \citep{hollander2010}.

Reading this language matters because the information in an undisclosed acquisition is exactly what securities law keeps off the call, yet the unscripted question-and-answer session is one of the few recurring venues in which a constrained executive must keep speaking and fielding questions. If the bind leaves a trace in the words, that trace would mark a deal's passage from private to public in the firm's own voice, in a channel separate from the price run-up already documented around mergers \citep{keown1981}.

The nearest prior work reads deal-related language along other dimensions. \citet{thewissen2024} show that acquirers manage the tone of their disclosure before stock-for-stock deals, and \citet{ragozzino2024} find that the volume of corporate-strategy vocabulary rises on the calls of acquisitive firms; both concern tone or subject matter rather than expressed uncertainty. Separately, \citet{keown1981} document that much of a merger's price reaction precedes its public announcement, an anticipatory signal that nonetheless lives in prices rather than in language. To our knowledge no prior work reads uncertainty language in the unscripted question-and-answer session in the anticipatory window before a withheld deal; we offer this as a positioning claim about where the contribution sits, not as a tested mechanism.

The purpose of this study is to ask whether chief-executive question-and-answer uncertainty tracks a deal's passage from private to public: whether it is elevated while an acquisition is withheld and recedes once it is announced. We construct a residual measure of chief-executive question-and-answer uncertainty, $\mathrm{UncResCEO}$, the call-varying component that remains once each executive's persistent speaking style is netted out, following a decomposition developed in the earnings-call language literature, and we defer its construction and the estimating equations to Section~2. We study how this measure behaves around cash acquisitions, using cash deals against a stock-acquisition comparison that shares the same withholding setting, on United States public, non-financial, non-utility firms' earnings calls from 2002 to 2018.

Our main findings can be previewed as follows. First, residual chief-executive uncertainty is elevated in the quarter before a cash acquisition, relative to the same firm's other quarters, with no comparable rise before stock acquisitions. Second, this elevation is anticipatory: it becomes indistinguishable from zero as soon as the deal is announced, even before it closes, whereas the cash that funds the purchase persists on the balance sheet until completion, two paths running on different clocks. This round-trip contrast is the study's strongest result.

Finally, the run-up concentrates in cash acquisitions rather than stock acquisitions, a difference that survives a formal pooled test of cash-specificity, though we read it as concentration rather than strict specificity, and the cash-accumulation mechanism behind it is left open. We also rule out the most immediate alternative. The run-up is not accounted for by analysts devoting more of the call to cash questions: with a direct measure of analyst cash scrutiny and its interaction entered alongside the pre-announcement indicator, the elevation survives and scrutiny does not amplify it. This means that scrutiny does not account for this particular run-up, not that scrutiny never matters. A further analysis turns to the firm's post-call information environment, measured by the bid-ask spread: the residual is unrelated to it, while the scripted presentation is positively associated with it, a pattern consistent with the scripted and unscripted segments reaching the outside environment in different ways.

This study makes four contributions, each descriptive. First, it reads a residual measure of chief-executive question-and-answer uncertainty, the call-varying component that remains once each executive's persistent speaking style is netted out, in the anticipatory window before an acquisition is public, where the nearest prior work has read managed tone, the volume of strategy vocabulary, and prices rather than uncertainty language. Second, it documents that this language pattern concentrates in cash acquisitions rather than stock, a contrast that survives a formal pooled test. Third, it reads the unscripted question-and-answer session as tracking a deal's passage from private to public, connecting the literature on when managers withhold private information with the evidence that markets anticipate corporate transactions, while claiming no causal channel. Fourth, it documents that this residual uncertainty is unrelated to the firm's post-call bid-ask spread while the scripted presentation is positively associated with it, consistent with a difference between the scripted and unscripted segments in how they relate to the outside information environment.

Throughout, we read these patterns as correlational and within-firm: the design supports no causal identification and establishes no mechanism. The concentration in cash deals is consistent with stock acquirers' incentive to manage their pre-deal narrative, an incentive that cash acquirers lack and that we develop in Section~2; we offer this as motivation, not identification. The source of the uncertainty, a compliance-constrained inability to speak versus a strategically chosen reticence, remains observationally equivalent.

The remainder of the paper proceeds as follows. Section~2 develops the conceptual framework and empirical strategy; Section~3 presents the three main analyses, the pre-announcement run-up, the differential timing around announcement, and the cash-specificity test; Section~4 reports additional analyses, ruling out analyst scrutiny, examining outsider reactions through the bid-ask spread, and a set of robustness checks---treating a withdrawal as a resolution event, dropping the cash equation's dynamic term, and re-running the main findings without the first-deal restriction; and Section~5 concludes.

\section{Theoretical Framework and Hypothesis Development}

\subsection{2.1}

Consider a firm that has privately committed to an acquisition but has not yet announced it. Such a firm sits in a particular disclosure state, not a particular balance-sheet state: the pending deal is material information that the public does not yet hold, while the quarterly earnings call is a standing obligation the firm must still meet. Disclosure theory gives this bind its shape. \citet{verrecchia1983} shows that a manager who knows something may rationally choose not to reveal it when disclosure is costly; there is a threshold below which the informed manager simply withholds. \citet{dye1985} shows why such silence can survive in equilibrium: outside investors cannot always tell a manager who knows nothing from one who knows something and is keeping quiet, so withholding is not, by itself, a giveaway. The pending deal can matter to investors well before any contract is signed. In \citet{basic_v_levinson}, the Court held that information about preliminary, still-pending merger negotiations can be material, and that materiality turns on weighing how likely the deal is against how large it would be -- a probability-times-magnitude assessment rather than a bright line. Yet the same law that makes the deal material also limits what the manager may say about it. A firm has no general duty to disclose confidential merger talks, so staying silent is permitted; but once it speaks, it may not mislead, because an untrue or half-true statement is unlawful \citep{basic_v_levinson, rule_10b5}. Hence a flat denial is not a lawful option, which leaves deflection. The result is a tight bind: the chief executive must host the call, cannot address the one material thing, and cannot deny it either. That bind -- a state of informed withholding, which the theory treats as itself informative -- is the organizing primitive of this paper.

The earnings call is where this bind becomes visible. It is a recurring, semi-spontaneous event, and its question-and-answer (Q\&A) session is far less scripted than the prepared remarks that open it. A growing body of work reads the language of these calls, not only their numbers. \citet{matsumoto2011} find that the analyst-discussion session carries information beyond the managers' prepared presentation -- the unscripted exchange adds something the script does not. Reading that language quantitatively is, by now, an established method, and it is the method this paper adopts rather than invents. Off-the-shelf word lists built for ordinary English misclassify financial text; \citet{lm2011} build finance-specific word lists -- among them the uncertainty list our measure uses -- to classify tone and uncertainty in financial disclosures. We take this apparatus as given. Our contribution is not the measurement but the use we put it to, which we set out below.

Now place a constrained chief executive in that Q\&A. If the one pressing question is the one he may not answer, his unscripted replies have nowhere clean to go: they may turn hedged, qualified, and non-committal, and the share of uncertainty words in his answers may rise. This is not a claim that managers speak at random. \citet{hollander2010} show that managers make deliberate disclosure-and-silence choices during conference calls, and that a non-answer or a silence is itself informative to the market -- saying little is a choice the market reads. To the extent that such a signal is keyed to the undisclosed deal, it should be present while the deal is still secret and should resolve once the announcement makes the deal public. The undisclosed state, in other words, supplies a clock: the language moves with the information and settles when the information comes out.

One caution is needed before going further. Some of the uncertainty in how a manager speaks is simply persistent personal style. \citet{bertrand_schoar2003} show that managers carry durable, individual styles that show up as manager fixed effects in firm policies and outcomes -- a measurable component that stays with the person, not the moment. Call language is no exception. \citet{dwz2021} separate a manager's uncertainty and clarity language in earnings calls into two pieces: a persistent, manager-specific component -- a chief-executive fixed effect -- and a time-varying, call-level residual that changes from one call to the next. This paper follows that split. The distinction matters here for one reason. An anticipatory signal that tracks a particular pending deal cannot be a fixed trait, because it must appear and then disappear within a single executive's tenure. Hence, to the extent that such a signal exists at all, it can only live in the call-varying part, once persistent style has been netted out. It is worth adding what that residual does not do elsewhere: \citet{dwz2021} find that the time-varying, residual component of call uncertainty is largely unrelated to market and stock-price reactions, so on its own it explains little for prices. An anticipatory, deal-tracking reading of the residual is therefore a claim about where a signal would have to sit, not a claim that the residual is already known to carry one.

Cash deals and stock deals share the same disclosure bind, but they differ in two ways that matter for the design. The first difference is on the balance sheet: a cash purchase draws on an accumulated, visible cash position -- a war chest, in effect -- that a stock deal need not assemble. \citet{harford1999} documents that firms build up cash reserves ahead of acquisitions and that cash-rich firms are more acquisitive, so a cash bid reflects an accumulated balance-sheet position. The second difference is about incentives to manage the story beforehand. A stock bidder pays in its own equity -- its acquisition currency -- and so has reason to care what that equity is worth. \citet{shleifer_vishny2003} argue that an overvalued bidder has a powerful incentive to keep its equity overvalued so that it can buy with stock, a valuation-and-currency motive that a cash bidder does not share. Bidders act on that motive before the deal: \citet{louis2004} finds that stock-for-stock bidders overstate reported earnings -- positive abnormal accruals -- in the quarter just before the announcement, whereas cash acquirers do not. The same pre-deal management reaches disclosure tone as well, since \citet{thewissen2024} report that stock bidders inflate the tone of their earnings press releases before a stock-for-stock acquisition. Two clarifications keep this in its proper place. First, on what the comparison buys us: because a cash bidder has no equity currency to protect, it lacks the parallel pre-deal management incentive that a stock bidder has, so we read the cash deal as the relatively unmanaged window onto the disclosure state and the stock deal as the managed comparison -- the same bind, but overlaid with an extra reason to manage the message. This is what motivates our focus on cash deals; how large the cash-versus-stock gap is, and why the run-up should concentrate in cash, are developed with H1a in the hypotheses that follow. Here the point is motivation, not identification, and the source of the cash uncertainty stays open. Second, a firewall between channels: the management documented above operates on scripted, prepared artifacts -- the tone of an earnings press release, the accruals in a reported number -- whereas our dependent variable is uncertainty in the unscripted Q\&A, the call segment managers cannot fully prepare in advance. Because we read the venue that is hardest to stage-manage, the scripted-channel management those papers describe does not mechanically carry into our measure; our signal is therefore not an artifact of the very tone-management the citations describe, which is a strength of the design rather than a gap in it. The two differences then come apart in time: uncertainty tracks the information and resolves at the announcement, while the cash position serves the purchase and, mechanically, persists until the deal completes.

Where does this leave our paper relative to the nearest work? The closest study reads how deal-related language is managed. \citet{thewissen2024} is nearest to us in setting -- stock-deal acquirers, before the announcement -- but its object is different: managed tone in scripted press releases, whereas ours is unmanaged uncertainty in the unscripted Q\&A. That difference places our signal in a distinct cell of the deal-language space and forecloses reading it as merely managed deal tone. Adjacent work reads other features of the same talk; \citet{ragozzino2024} study the volume of strategy vocabulary on calls around deal activity. A separate literature documents leakage on the price side: \citet{keown1981} show abnormal stock-price run-up before public merger announcements, so that pre-announcement information leaks into prices. What none of this reads is unmanaged language along an uncertainty dimension, before a cash deal, in the anticipatory window. To our knowledge, no prior work occupies that cell, and that empty cell is where this paper sits.

One last boundary. The prediction admits two readings, and we do not choose between them. On one reading the disclosure is compliance-constrained: the chief executive legally cannot speak to the deal. On the other it is strategic silence: the chief executive chooses not to. In our data these two are observationally equivalent -- the same pattern is consistent with either -- and the framework does not require us to separate them. We therefore document a pattern keyed to the disclosure state, claim no identification, and leave the underlying mechanism open.

\subsection{2.2}

We now turn the framework into predictions we can test. The two dimensions developed above -- anticipatory timing and cash-concentration -- speak to where and when a chief-executive Q\&A language signal should appear, not to why it appears. To fix the timing, we define the pre-announcement quarter. Let event time $e$ be the call quarter minus the announcement quarter, so that $e = -1$ is the call held in the quarter immediately before the firm's first acquisition that is at least 50\% cash. We then set $\mathit{PreAnnounceQtr} = \mathbf{1}[e = -1]$, an indicator equal to one for that single pre-announcement call and zero otherwise. We state the predictions as falsifiable hypotheses -- H1, H1a, and H1b -- and test them with the designs that follow. Throughout, we read the results in descriptive, correlational terms: they describe where and when the signal sits, and we make no causal claim.

Our first prediction concerns the timing of the signal for cash acquirers. \begin{quote} \textbf{H1.} Residual chief-executive Q\&A uncertainty (\textit{UncResCEO}) is elevated in the pre-announcement quarter of cash acquirers, relative to those same firms' other quarters. \end{quote} Here \textit{UncResCEO} is the call-level residual uncertainty defined in the framework -- the part of a manager's answer-uncertainty that is left once persistent personal style is removed. H1 is not the obvious prediction, because prior work documents the opposite register. \citet{thewissen2024} show that stock bidders deliberately manage their disclosure tone before stock-for-stock deals, dressing the message up on purpose. H1 predicts something different in kind: not managed tone before stock deals, but unmanaged uncertainty before cash deals. What makes H1 non-obvious is exactly this contrast -- the literature describes deliberate, strategic tone management, whereas H1 predicts that uncertainty language surfaces under the disclosure bind. Consistent with the boundary drawn above, H1 takes no stance on whether that surfacing is strategic or compliance-constrained; it predicts that the language appears, not why.

Our second prediction sharpens the first by comparing cash deals with stock deals. \begin{quote} \textbf{H1a.} The pre-announcement uncertainty run-up is stronger for cash acquirers than for stock acquirers. \end{quote} The reasoning runs through the two differences set out earlier. A cash bid draws on an accumulated cash position, while a stock exchange need not; \citet{harford1999} documents the cash build-up that precedes acquisitions. Moreover, a stock bidder carries a currency-protection motive that a cash bidder lacks. This is why we treat the stock deal as a managed comparison rather than an inert placebo: it faces the same disclosure bind, but with an additional incentive to manage the message beforehand. Two careful limits apply to H1a. First, it is a prediction about the effect -- the run-up concentrating in cash deals -- and we frame it as concentration, not strict specificity; H1a does not claim that the cash build-up itself is cash-specific, and it asserts no cause. Second, because the stock arm is a managed comparison, any pre-deal signal there is entangled with documented management -- a competing story that the cash setting does not carry. That asymmetry is what motivates the predicted direction, the run-up concentrating in cash. It motivates the direction of the prediction; it does not detect anything. We interpret, and we do not detect.

Our third prediction is about timing after the announcement, and it turns on two different clocks. \begin{quote} \textbf{H1b.} Residual Q\&A uncertainty resolves at the announcement -- falling even before the deal closes -- while the acquirer's cash persists until it is paid at completion. \end{quote} Uncertainty runs on an information clock: once the announcement makes the deal public, the residual should drop, including in the gap window -- the stretch when the deal is announced but not yet completed. Cash runs on a transaction clock: it stays on the balance sheet until it physically leaves at closing. Two cautions attach to H1b. The cash-persistence leg is partly mechanical, since the cash does not leave until the deal closes. And the post-announcement undershoot in uncertainty should not be over-read; we flag the timing contrast, not a precise magnitude.

One alternative reading deserves to be flagged here, though we develop and test it later. The competing hypothesis is analyst scrutiny: perhaps the pre-announcement rise in uncertainty reflects analysts asking more cash-related questions, which mechanically draws out hedged answers, rather than the disclosure bind doing the work. This is an identification threat. We motivate and develop it where we set out our construct measures, and we test it as a side analysis with the results, where the run-up survives it -- though we read that as an underpowered failure to find, not a powered equivalence test. Two points keep this in proportion. It is distinct from the two observationally-equivalent readings drawn earlier -- compliance-constrained versus strategic silence -- which concern why the manager is quiet, not whether scrutiny is doing the work. And ruling this channel out does not establish a mechanism; it removes one rival story, nothing more.

\subsection{2.3}

We measure uncertainty with the finance-specific apparatus introduced earlier. Following \citet{lm2011}, we take their `uncertainty' word list -- a list built for financial text -- and apply it as a word-share, the fraction of words that fall on the list. This is the established tool, not a new one. On top of it we build the raw measure that the decomposition operates on. Following \citet{dwz}, we define \[ \mathit{UncAnsCEO} = \frac{\mathit{UnctWordsAnsCEO}}{\mathit{WordsAnsCEO}}, \] that is, the number of uncertainty words the chief executive uses when answering analysts divided by the total number of words spoken in those answers (their equation~(2)). In words, it is the share of the chief executive's answer that is uncertainty language. This raw share is the input to the decomposition and the left-hand side of the decomposition equation below. We compute the same word-share on the two parts of the call separately: \textit{UncPreCEO} is the share in the chief executive's prepared presentation (their equation~(1)), and \textit{UncQue} is the share in the analysts' questions (their equation~(3)). Both take the same form as \textit{UncAnsCEO}, but they enter the decomposition as controls rather than as the dependent variable.

Why a residual at all? As established earlier, an anticipatory signal that tracks a particular pending deal must vary from one call to the next within a single chief executive's tenure, so it cannot live in the persistent, manager-specific style; it can only live in the call-varying residual. The decomposition isolates exactly that residual. We use the specification of \citet{dwz}, estimated on our own data: \[ \mathit{UncAnsCEO}_{it} = \alpha + \gamma_i\,\mathit{CEO}_{it} + \beta_s\,\mathit{Speech}_{sit} + \beta_k\,\mathit{FirmChars}_{kit} + \mathit{Year}_t + \varepsilon_{it}, \] where $\alpha$ is a constant; $\mathit{CEO}_{it}$ picks out the chief executive, so that $\gamma_i$ is that manager's persistent style and $\mathit{ClarityCEO}_i = -\gamma_i$; $\mathit{Speech}_{sit}$ collects the other call-language controls $\{\mathit{UncPreCEO}, \mathit{UncQue}, \mathit{NegCall}\}$; $\mathit{FirmChars}_{kit}$ collects the firm controls $\{\mathit{SurpDec}, \mathit{EPSgrowth}, \mathit{StockRet}, \mathit{MarketRet}\}$; and $\mathit{Year}_t$ is a set of year fixed effects. The object we carry forward is the residual, $\mathit{UncResCEO} = \varepsilon_{it}$, the part of answer-uncertainty left unexplained by all of these. One point about ownership is worth stating plainly: this is the specification of \citet{dwz}, re-estimated on our sample, so the residual is our own, not their published residual. This equation is the precise version of the `call-level residual' described earlier. It nets out two things at once: the persistent style $\gamma_i$, and observable call-level factors -- the prepared-presentation uncertainty, the analysts' question uncertainty, call negativity, contemporaneous performance, and the year. Because it removes more, \textit{UncResCEO} is narrower than a style-only residual would be; it is a refinement of the earlier description, not a different object. The measure itself is not new to us: \citet{dwz} construct \textit{UncResCEO} and use it as an explanatory variable in their own model of market responses (their equation~(5)). Our contribution is not the measure but the reading we give it -- an anticipatory, disclosure-state reading of \textit{UncResCEO} around deals that have not yet been disclosed. Following \citet{dwz}, the residual is answer-uncertainty net of the scripted presentation (\textit{UncPreCEO}): it isolates the unscripted, call-specific component, which is where the anticipation of an undisclosed deal would surface. If some of that anticipation also leaks into the scripted, vetted remarks, then the residual understates it, which makes the measure a conservative floor rather than an upper bound. By the same construction the residual is also net of the analysts' own question-uncertainty (\textit{UncQue}) and of call negativity (\textit{NegCall}), and we rely on this when we reconcile the measure with the scrutiny and tone alternatives later.

One property of this construction deserves a note, offered as a caution rather than a complaint. \textit{UncResCEO} is a generated regressand -- a residual estimated in a first step that we then use as the dependent variable in a second step. \citet{pagan1984} sets out the issue for such two-step settings. This is a standard arrangement, not a flaw peculiar to us; indeed \citet{dwz} themselves use the residual as a regressor in their own second step. The practical implication is that conventional standard errors may be understated unless a two-step correction is applied. We flag the concern, together with the bootstrap that would address it, and note that it does not change the descriptive readings we report. A second property we do not assume away: whether the residual truly captures uncertainty, rather than merely reflecting how hard analysts are pressing, is an empirical question. We treat it as one and take it up where we examine the construct's validity, rather than asserting it here.

\subsection{2.4}

We estimate three related designs on a common panel. The first and main one, which we label MA1, is \[ Y_{it} = \beta\,\mathit{PreAnnounceQtr}_{it} + \gamma' X_{it} + \alpha_i + \tau_t + \varepsilon_{it}, \] where $Y_{it}$ is the outcome for firm $i$ in quarter $t$; $X_{it}$ is a vector of firm controls $\{\mathit{Leverage}, \mathit{lnAssets}, \mathit{TobinsQ}, \mathit{ROA}, \mathit{Capex}, \mathit{DivDummy}, \mathit{sCFO}\}$; $\alpha_i$ is a firm fixed effect, which absorbs everything that stays fixed about the firm; and $\tau_t$ is a calendar year-quarter fixed effect, which absorbs whatever is common to all firms in that quarter. We cluster standard errors by firm. The treatment indicator is $\mathit{PreAnnounceQtr}_{it} = \mathbf{1}[e = -1]$, where event time $e$ is the call quarter minus the quarter of the firm's first acquisition that is at least 50\% cash. We drop the post-announcement quarters ($e \geq 0$) from the sample, and firms that never acquire (for which $e$ is undefined) serve as the fixed-effects baseline. The coefficient $\beta$ is the estimand: it measures a within-firm shift in the outcome in the pre-announcement quarter -- the anticipation quarter -- read against the same firm's own earlier, normal quarters. The never-acquirers mainly pin down the calendar-time effects rather than serving as the within-firm contrast. We read $\beta$ as a descriptive contrast, not a causal effect.

The second design, which we label MA2, replaces the single indicator with a small set of event-time bins, so that we can see the shape over time: \[ Y_{it} = \sum_b \beta_b\,\mathit{Bin}_{b,it} + \gamma' X_{it} + \alpha_i + \tau_t + \varepsilon_{it}. \] There are four bins. \textit{PRE2} is two quarters before ($e = -2$); \textit{PRE1} is the pre-announcement quarter ($e = -1$); \textit{GAP} covers the quarters after the announcement when the deal is announced but not yet closed; and \textit{POST} covers the quarters after the deal has completed. The omitted baseline is the set of quarters three or more before the announcement ($e \leq -3$) together with the never-acquirers. We cap the post-window at four quarters, truncate it at a firm's second announcement, and drop the post rows of deals that are later withdrawn. The discriminating contrast is a Wald test on $\beta_{\mathit{PRE1}} - \beta_{\mathit{GAP}}$: if uncertainty drops at the GAP -- announced but not yet completed -- that favours a reading in which the constraint lifts at announcement over one in which the outcome only resolves at completion. \textit{PRE2} is the validity check; with no anticipation that early, its coefficient is expected to be about zero, which is what a clean pre-trend would show. We run the event study for cash acquirers and, as a benchmark for comparison, for stock acquirers, and we report all bin coefficients and Wald drops two-tailed.

The third design, which we label MA3, puts both treatments in one model so that we can compare them directly: \[ Y_{it} = \beta_c\,\mathit{PreAnn}^{\mathrm{cash}}_{it} + \beta_s\,\mathit{PreAnn}^{\mathrm{stock}}_{it} + \gamma' X_{it} + \alpha_i + \tau_t + \varepsilon_{it}, \] where $\mathit{PreAnn}^{\mathrm{cash}}_{it} = \mathbf{1}[cq = dq_{\mathrm{cash}} - 1]$ marks the quarter just before a cash deal and $\mathit{PreAnn}^{\mathrm{stock}}_{it} = \mathbf{1}[cq = dq_{\mathrm{stock}} - 1]$ marks the quarter just before a stock deal, with $cq$ the call quarter and $dq$ the deal quarter. The sample keeps only the quarters before a firm's first deal of either type, plus firms that do neither; those never-either firms are the fixed-effects baseline. The formal test of cash-specificity is a Wald test of whether $\beta_c - \beta_s > 0$. This is a single linear-restriction test on the difference, not the weaker move of comparing one significant coefficient against one that is not. One specification detail applies only to the cash-position outcome: for the \textit{CashRatio} outcome -- the regression in which the firm's cash-to-assets position is itself the dependent variable -- we add that variable's own one-quarter within-firm lag, a partial-adjustment term using consecutive quarters only, so that the pre-announcement coefficient tracks the change in cash. The residual \textit{UncResCEO} outcome takes no lagged dependent variable.

Three assumptions are common to all three designs, and we state them descriptively. Identification is within-firm, through the firm fixed effects, and within calendar-time, through the year-quarter fixed effects. The estimand is the pre-announcement mean shift, read against the firm's own normal quarters and against non-acquirers. The \textit{PRE2} bin is the validity device that lets a reader judge the pre-trend. We make no parallel-trends assumption and no causal claim. The three designs map onto the three hypotheses one for one: MA1, with its single pre-announcement indicator, tests H1; MA2, with its event-time bins, tests H1b; and MA3, with its pooled cash-versus-stock interaction and Wald test, tests H1a. Two distinct control sets are also at work, and it helps to keep them apart. The first-stage residualization controls -- the call-language and performance variables, with chief-executive and year fixed effects -- are the ones that build \textit{UncResCEO}. The second-stage design controls -- the firm financials, with firm and year-quarter fixed effects -- are the ones in the equations above. Because the outcome \textit{UncResCEO} is already net of the first-stage set, the second-stage designs simply add the firm-financial controls and a finer time fixed effect on top.

Several threats do not belong to any single equation, and we disclose them together. First, inference: we cluster standard errors by firm in the three designs, and some validity tables use two-way clustering, by firm and by calendar quarter. In places we report the focal treatment tests one-tailed; for honesty these should be read two-tailed, and the load-bearing results survive that stricter standard -- in particular, the pre-announcement run-up remains significant at $p = 0.0074$ two-tailed. Second, the endogeneity of timing: firms choose when to do deals, and only the first deal sets our event clock, so later deals can contaminate the baseline; we mitigate this by truncating each firm's post-window at its second announcement, but the never-acquirer baseline and the focus on the first deal remain assumptions, not identification. Third, functional form: the outcomes are levels and word-shares, not counts; the \textit{CashRatio} regressions use a partial-adjustment lag, because cash is sticky -- its quarter-to-quarter persistence is about $\rho \approx 0.78$ -- while the residual outcomes use no such lag; and because \textit{UncResCEO} is a generated regressand, the two-step standard-error caveat noted earlier applies to every design here. Fourth, the choice of controls: the firm-financial controls -- \textit{Leverage}, \textit{lnAssets}, \textit{TobinsQ}, \textit{ROA}, \textit{Capex}, \textit{DivDummy}, and \textit{sCFO} -- follow the standard cash-determinants regression \citep{opler1999, bates2009}, and we ground the cash equation in that literature and read it mechanically, as cash accumulating as a by-product of retained free cash flow rather than as purposeful saving toward a war chest. Fifth, sample selection. The residual outcome \textit{UncResCEO} exists only for chief executives with at least five conference calls in the main, non-financial and non-utility, industry sample -- the universe in which the first-stage residual is estimated. Because every design drops observations that lack the dependent variable, all of the focal uncertainty regressions for H1, H1b, and H1a run on this subsample of firms with at least five calls in the main industries, which skews toward larger, better-covered firms. The \textit{CashRatio} outcomes, by contrast, use the full panel. We disclose this as a selection threat to external validity.

\subsection{2.5}

Having defined the residual, we now check whether it measures what we claim, and we treat that validity as an empirical question to answer here, not as an assumption. We ask two things of $\mathrm{UncResCEO}$. First, is it convergent: does it move together with measures of uncertainty that are already established? That is the check we run in this section. Second, could a plausible rival account for the pre-announcement pattern instead of the disclosure bind, namely analyst scrutiny, how much of a call analysts spend asking about cash? That rival is a plausible alternative driver, and we put it to the data as a side analysis in Section~4.1. By construction, the residual is already net of the scripted presentation and of the wording of the analysts' own questions, as set out in Section~2.3, so these checks ask how it behaves once that wording has been removed.

We begin with convergent validity. The residual is positively associated with three established measures of uncertainty. It rises with firm-level political risk (PRisk), with a coefficient of $0.0001^{***}$; with US economic policy uncertainty (US-EPU), at $0.0124^{**}$ and $0.0123^{*}$, the second estimate only marginally significant once firm fixed effects are included; and with global economic policy uncertainty (GEPU), at $0.0181^{**}$ and $0.0187^{**}$. We read these as consistent with established measures, not as establishing the construct, and three cautions keep that reading honest. First, the tests are one-tailed. Second, although political risk is measured for each firm each quarter, exactly the same unit as the residual and so the cleanest match of the three, its association is economically trivial, with an $R^2 \approx 0.003$. Third, US-EPU and GEPU are macro, monthly indices: in any given period every firm shares the same value. Because these enter under calendar-year fixed effects, they are identified only by within-year aggregate co-movement, that is, only by whether the residual and these economy-wide indices move together within a year. That is an aggregate co-movement check, not a firm-level one. In short, the convergent leg is supportive but weak.

We now turn to that rival, and we build our own measure to confront it. We define $\mathrm{CashScrutiny}$ as the share, in percent, of analyst question-and-answer turns devoted to cash and liquidity, and $\mathrm{HighCashScrutiny}$ as an indicator equal to one when $\mathrm{CashScrutiny}$ is above its median. The median is zero: only about 11\% of calls draw any cash scrutiny at all. This is a measure of volume, or topic: how much of the conversation is about cash. It is therefore distinct from the uncertainty wording of the questions ($\mathrm{UncQue}$). We motivate the side-test on two grounds. First, the measure is valid: the firm's cash ratio predicts it, with coefficients of $0.7530^{***}$ and $0.8519^{***}$, the same class of validity check as the established measures above. Second, it is a plausible alternative driver: cash scrutiny itself rises around cash-deal announcements, so it co-moves with both the firm's cash and the event, and it could therefore generate the pre-announcement run-up mechanically, which is a textbook identification threat. We therefore test, as a side analysis in Section~4.1, whether scrutiny rather than the disclosure bind accounts for the run-up. We do not preview the magnitudes of that test or its verdict here; the formal side-test, its hedged verdict, and the underpowered caveat are all reported in Section~4.1. One might object that this test merely repeats the controls already inside the residual. It does not. $\mathrm{CashScrutiny}$ measures the volume and topic of cash questions, whereas the residual already removes the uncertainty wording of those questions, $\mathrm{UncQue}$, by construction, through the decomposition of Section~2.3. So the test targets the question-volume path that the residual does not absorb, and it is not redundant with the controls already built into the residual.

Finally, we collect the remaining key constructs used in the designs. $\mathrm{CashRatio}$ is the firm's cash-to-assets position, the cash it holds relative to its assets. $\mathrm{PreAnnounceQtr}=\mathbf{1}[e=-1]$ is the pre-announcement-quarter indicator defined in Section~2.2 and used in Section~2.4. $\mathrm{CashScrutiny}$ and $\mathrm{HighCashScrutiny}$ are the analyst-attention measures defined just above. The remaining controls, $\mathrm{Leverage}$, $\mathrm{lnAssets}$, $\mathrm{TobinsQ}$, $\mathrm{ROA}$, $\mathrm{Capex}$, $\mathrm{DivDummy}$, $\mathrm{sCFO}$, and the controls from the Section~2.3 decomposition that builds the residual, are catalogued in the Appendix.

\section{Data and Main Results}

\subsection{3.1}

We combine five data sources. First, we download Capital~IQ earnings-call transcripts for U.S. public firms over 2002--2018, parsed so that we know each speaker's role and whether a turn falls in the scripted presentation or in the back-and-forth question-and-answer (Q\&A) segment. Second, we retrieve quarterly Compustat North America fundamentals, taking the most recent fiscal quarter that ends on or before the call. Third, we pull CRSP daily records and IBES earnings records, which supply the stock-return, market-return, and earnings-surprise controls that the Section~2.3 decomposition requires. Fourth, we collect SDC merger-and-acquisition records---announcement date, completion or withdrawal date, and the cash-versus-stock split of how each deal is paid for---keeping only deals with a transaction value of at least \$1~million. Fifth, we take Execucomp annual CEO records and aggregate them into a monthly tenure panel that identifies which executive holds the CEO role in each quarter; this identity is the input to the Section~2.3 manager fixed effect. These five sources do not all use the same firm identifier. SDC keys each deal by the acquirer's six-digit CUSIP, while our call panel is keyed by the Compustat firm identifier (gvkey), so we join deals to calls through the acquirer's six-digit CUSIP, using the CRSP--Compustat link table that anchors the master-sample manifest.

We build the language measures in-house. We first tokenize each call---that is, we break it into individual words---and then compute the Loughran--McDonald uncertainty-word share, the fraction of words drawn from a fixed list of uncertainty terms, separately for each speaker and each segment. The CEO's uncertainty share in the Q\&A, UncAnsCEO, is the input that the Section~2.3 first stage (the DWZ decomposition, estimated on our own sample) turns into the residual measure UncResCEO---the part of the CEO's answer uncertainty that is left after its predictable component is removed. This is the construction step only; the equation itself, and the definition and justification of UncResCEO, live in Section~2.3, and we do not re-derive them here.

The remaining key variables are constructed here but defined in full elsewhere, and we point to each home rather than re-derive it. The residual uncertainty measure UncResCEO is defined in Section~2.3. The pre-announcement indicator PreAnnounceQtr equals one in the single quarter just before a firm's first qualifying deal ($e=-1$) and zero otherwise; it is defined in Sections~2.2 and~2.4. The cash ratio, CashRatio, is cash and equivalents divided by total assets (\texttt{cheq}/\texttt{atq}), defined in Section~2.5 and the Appendix. Cash scrutiny, CashScrutiny, is the share of analyst Q\&A turns that touch on cash or liquidity, counted from the locked word list in Appendix~I; because the median call draws none at all (median $0.0000$), we also use a simple indicator, HighCashScrutiny, that flags any call with cash scrutiny present---only about 11\% of calls clear that bar (mean $0.1127$). Both are defined in Section~2.5. Finally, the firm-financial controls---leverage, log assets, Tobin's~Q, return on assets, capital expenditure, a dividend indicator, and cash-flow volatility---are catalogued in the Appendix.

The estimation samples are built in three layers, which is why the number of observations changes from one table to the next. The first layer is the base call panel. We keep every matched Capital~IQ call of a U.S. public firm that is neither a financial nor a utility company over 2002--2018. We drop any CEO with fewer than five calls, because with so few calls we cannot estimate a speaking style and therefore cannot form the residual. We also restrict the panel to firms that Execucomp covers---roughly the S\&P~1500---because the CEO's identity feeds the Section~2.3 manager fixed effect. The second layer is the SDC deal universe. We keep U.S. public-acquirer deals over 2002--2018 whose status is completed, pending, or withdrawn and whose payment composition is known. Deals that are at least 50\% cash form the cash arm; deals that are at least 50\% stock form the stock comparison arm; and each firm's first qualifying deal sets its event clock. The third layer is a set of design-specific trims. The run-up test drops the post-announcement quarters of treated firms. The event studies cap the post-window at four quarters and truncate it at the firm's next announcement. The scrutiny validity and channel tests require at least three analyst Q\&A turns per call. We state the five-call filter here purely as a construction mechanic; we take up its implications for sample selection in Section~2.4.

Table~\ref{tab:summary_stats} reports summary statistics over 2002--2018. Our earnings-call (language) sample comprises 88,205 earnings-call observations across 1,884 firms; the cash-ratio regressions draw on a broader panel spanning 2,232 firms, because the cash ratio does not require a CEO speaking-style residual and so is not limited to the calls that can form one. The count $N$ varies from row to row, because each variable is summarized only on the calls for which it is available (its complete cases). The residual uncertainty measure UncResCEO is centered near zero, with mean $-0.0000$ and a standard deviation of $0.3010$ (44,900 observations, Panel~B), while the cash ratio averages $0.1708$ with a standard deviation of $0.1806$ (Panel~C). Because the residual can be computed only for CEOs with at least five calls at non-financial and non-utility firms, the uncertainty regressions tend to describe a better-covered set of firms, one that skews toward larger firms, than the cash-ratio regressions do. For that reason every table reports its own firm and observation counts, and we are cautious about reading too much into these summary statistics on their own.

\subsection{3.2}

Our first main analysis (MA1) estimates the single-indicator two-way fixed-effects regression of Section~2.4 four times. The regression carries firm fixed effects and calendar year-quarter fixed effects---in plain terms, it compares a firm against itself and against the same calendar quarter at other firms---and clusters standard errors by firm. We run it once for each outcome---the cash ratio (CashRatio) and the residual uncertainty in the CEO's Q\&A answers (UncResCEO)---in each arm: the cash arm, deals that are at least half cash, and the stock comparison arm, deals that are at least half stock. For treated firms we drop the quarters after the announcement, so the design sees only the run-up into the deal. Identification is within-firm: it sets a treated firm's pre-announcement quarter against that same firm's ordinary quarters. This is a correlational design and does not by itself identify a causal effect; the estimating equation is given in Section~2.4, and we do not restate it here.

Table~\ref{tab:empire_building_did} shows that in the cash arm both measures move in the final quarter before the announcement. Column~2 reports that residual Q\&A uncertainty is elevated, with a coefficient of $0.0461^{***}$ (standard error $0.0172$)---the CEO's answers carry more uncertainty in that pre-deal quarter. Column~1 reports that the cash ratio is rising, with a coefficient of $0.0051^{***}$ (standard error $0.0017$); this coefficient should be read as a change in cash, not a level, because the cash equation carries a partial-adjustment lag of $0.7990^{***}$ (standard error $0.0070$)---that is, most of this quarter's cash is already explained by last quarter's, so the coefficient picks up only the increment added on top. One might worry that the uncertainty result is an artifact of the table's one-tailed reporting convention; it is not. The residual coefficient survives a two-tailed test ($p=.0074$), so the load-bearing result does not depend on the reporting convention. Because each firm is compared with itself, the pattern is within-firm; it remains correlational and does not identify a cause.

The two cash-arm columns are not run on the same set of firms. Because the residual can be formed only for CEOs with at least five calls, the uncertainty test in column~2 uses 27,622 firm-quarters from 1,248 firms, while the cash-ratio test in column~1 uses 67,590 firm-quarters from 2,232 firms. The cash-ratio panel therefore spans more firms (2,232) than the 1,884-firm residual-feasible call sample, because the cash ratio is available for firms whose CEOs lack the five calls the residual requires. The uncertainty test is therefore the more demanding one, run on the smaller universe.

How large are these effects? They are material but modest. The uncertainty coefficient of $0.0461^{***}$ is about fifteen percent of a residual standard deviation---roughly 15.3\%, taking $0.0461$ against the all-universe standard deviation of $0.3010$ reported in Table~\ref{tab:summary_stats}---and it appears in a measure from which everything persistent has already been removed, concentrated in a single quarter that makes up only about one percent of the sample (the pre-announcement indicator has a mean of $0.0094$). The cash coefficient of $0.0051^{***}$ adds roughly half a percentage point of assets to the cash stock---about three percent of the mean cash ratio ($0.0051$ against $0.1708$)---in just one quarter. These magnitudes are supportive but not definitive, and they do not establish a cause.

The stock arm is our comparison. It faces the same disclosure constraint, and its model is almost identical to the cash arm's---the stock partial-adjustment lag of $0.8013^{***}$ (standard error $0.0057$) is nearly the same as the cash arm's $0.7990^{***}$. Yet neither outcome shows a comparable positive run-up. In column~6 residual uncertainty carries a coefficient of $-0.0429$ (standard error $0.0307$), which is not statistically significant; in column~5 the cash ratio carries $-0.0036$, also not significant. We read these as flat, noisy nulls, not as movements in any direction. What differs at $e=-1$, the final quarter before the announcement, is therefore the response, not the model. We should be careful here: reading the two arms side by side is not yet a formal test of the difference between them. That formal cash-versus-stock test is supplied in Section~3.4, our third main analysis.

\subsection{3.3}

Our second main analysis (MA2) replaces the single pre-announcement indicator with a four-bin event study. The bins trace the deal through time: \textit{PRE2} marks the quarter two before the announcement ($e=-2$), \textit{PRE1} the quarter just before it ($e=-1$), \textit{GAP} the window in which the deal is announced but not yet closed, and \textit{POST} the quarters after completion. The baseline against which these are measured is the firm's own ordinary quarters ($e\le-3$) together with firms that never acquire. We estimate the specification on the matched universe---the set of firm-quarters where both the residual and the cash ratio are defined, 28,102 firm-quarters across 1,320 firms---and we run both outcomes on exactly these same rows. Because we impose no ordering across the bins, every coefficient is read two-tailed. Putting both outcomes on the identical sample guards the differential-timing reading against the worry that the two paths merely reflect different samples. The design is otherwise that of MA1, and it remains correlational: it does not identify a cause.

Table~\ref{tab:empire_drop_matched} first lets us check for a pre-trend---a drift in the outcomes that is already under way before the deal. There is none on either clock. The \textit{PRE2} coefficient is $0.0068$ (standard error $0.0178$) for the residual and $0.0008$ (standard error $0.0024$) for the cash ratio, neither significant. So the \textit{PRE1} estimates that follow are not simply the tail end of a trend that predates the announcement. We treat this as a validity check, not as a proof of identification.

The same table delivers the discriminating contrast: what happens once the deal becomes public, in the gap window between announcement and closing. On the information clock---the path traced by residual uncertainty, in column~1 of Table~\ref{tab:empire_drop_matched}---the residual spikes at \textit{PRE1} to $0.0473^{***}$ (standard error $0.0178$) and then falls to insignificance the instant the deal is announced, with a \textit{GAP} coefficient of $0.0018$ (standard error $0.0187$), not significant. The drop from \textit{PRE1} to \textit{GAP} is itself significant, $0.0455^{**}$ (standard error $0.0225$), and the full swing from the pre-announcement peak to completion is large and significant, $0.0723^{***}$ (standard error $0.0190$). On the transaction clock---the path traced by the cash ratio, in column~2---the picture differs. Cash is elevated at \textit{PRE1}, $0.0061^{**}$ (standard error $0.0024$), but it does \emph{not} fall when the deal is announced---the \textit{PRE1}-to-\textit{GAP} drop is just $0.0006$, not significant---and it declines only once the cash actually leaves the firm at completion, a \textit{GAP}-to-\textit{POST} drop of $0.0210^{***}$ (standard error $0.0042$) and a \textit{POST} level of $-0.0155^{***}$ (standard error $0.0022$). One caution is in order: the \textit{GAP} cash coefficient itself is $0.0055$ (standard error $0.0034$), positive but not significant, so the persistence of cash through the gap rests on the \emph{absence} of a \textit{PRE1}-to-\textit{GAP} decline, not on a significantly elevated gap level. Throughout, firms are compared with themselves, the design is correlational, and it does not identify a cause.

The entry spike is, in economic terms, the same event we saw in MA1: $0.0473^{***}$ here against $0.0461^{***}$ in Table~\ref{tab:empire_building_did}, about fifteen percent of a residual standard deviation (roughly $15.7\%$ on the $0.3010$ basis). What MA2 adds is the round trip. The swing from the pre-announcement peak to the post-completion quarters, $0.0723^{***}$, says that the uncertainty raised in the run-up is fully unwound once the information becomes public. The cash leg, meanwhile, uses the same well-behaved partial-adjustment specification as MA1, with a lagged-cash coefficient of $0.7547^{***}$ (standard error $0.0108$). This remains a correlational reading, not a causal one.

Column~1 of Table~\ref{tab:empire_drop_placebo} corroborates the round trip on a second sample. On its own complete-case cash universe of 29,535 firm-quarters from 1,326 firms, the cash arm of the comparison event study reproduces the same pattern: the residual again spikes at \textit{PRE1} ($0.0486^{***}$, standard error $0.0174$) and again collapses to insignificance at the gap ($0.0058$, standard error $0.0181$, not significant); the \textit{PRE1}-to-\textit{GAP} drop is significant ($0.0428^{*}$, standard error $0.0221$) and the peak-to-post swing is significant ($0.0681^{***}$, standard error $0.0187$). For now we use only the cash column of this comparison table; its cash-versus-stock contrast is reserved for Section~3.4.

Taken together, the claim rests on the contrast between the two paths, not on the size or sign of any single coefficient. Uncertainty tracks the flow of information and reverses the moment the deal is announced; cash tracks the transaction and falls only at completion---a persistence that is partly mechanical, since the cash does not leave the firm until the deal closes. We do not over-read the one coefficient that dips below baseline, the \textit{POST} uncertainty estimate of $-0.0250^{*}$ (standard error $0.0130$). The reading is correlational and supportive, not definitive; it does not identify a cause, and the mechanism behind it remains open.

\subsection{3.4}

We begin descriptively, rerunning the event study side by side on the stock comparison arm, using the stock column of Table~\ref{tab:empire_drop_placebo} that we held back from Section~3.3. Neither arm trends into the event: the PRE2 coefficient is flat in both ($0.0105$ in the cash arm and $-0.0056$ in the stock arm, neither significant). Only the cash arm then moves. It reproduces the run-up (PRE1 $0.0486^{***}$, standard error $0.0174$; peak-to-post drop $0.0681^{***}$, standard error $0.0187$), while the stock arm shows none of it (PRE1 $-0.0404$, standard error $0.0299$, not significant) and, if anything, its PRE1-to-GAP change is a marginal $-0.0756^{*}$ (standard error $0.0405$), a noisy across-announcement change we do not read as a directional run-up. This is a descriptive side-by-side, not yet the test.

A side-by-side display implies a comparison it never actually runs. So our third main analysis (MA3) puts both treatments into one pooled, interacted model --- the model of Section~2.4 --- and tests the cash-minus-stock difference directly, with a single joint test (a Wald linear restriction, $\beta_c-\beta_s>0$). That restriction is what estimates $\theta_{-1}^{\mathrm{cash}}-\theta_{-1}^{\mathrm{stock}}$ in hypothesis~H1a. This difference, and not a comparison of one significant coefficient against one insignificant one, is the formal cash-specificity test. The pooled baseline is the set of firms that make neither a cash nor a stock deal. The design remains correlational and identifies no causal effect, and it speaks to concentration in cash deals rather than strict specificity.

On the effect itself, column~1 of Table~\ref{tab:empire_cashspec} delivers the formal cash-specificity result. The cash arm's pre-announcement uncertainty coefficient is significant ($0.0459^{**}$, standard error $0.0185$), while the stock arm's is insignificant and negative ($-0.0524$, standard error $0.0436$). The formal cash-minus-stock difference, the Wald restriction, is $0.0983^{**}$ (standard error $0.0476$; $z\approx2.07$, $p=.039$), significant at the five-percent level two-tailed. Because the cash arm is positive while the stock arm is an imprecise null near zero, the gap between them is larger than either coefficient on its own --- roughly a third of a residual standard deviation (about 32.7\% of $0.3010$); the gap therefore rides in part on the imprecise negative stock estimate. We emphasize that this formal linear-restriction test, and not a comparison of one significant against one insignificant coefficient, is what we mean by the cash-specificity claim, and we read it as supportive rather than definitive.

On the proposed cause, the picture is weaker. The cash build-up itself is only faintly cash-sided: the cash-minus-stock difference in the cash ratio is insignificant on the matched universe (column~2, $0.0064$, standard error $0.0071$) and reaches only marginal significance on the full panel (column~3, $0.0092^{*}$, standard error $0.0055$). So while the \emph{effect} concentrates in cash deals, the proposed war-chest mechanism --- that firms deliberately build cash ahead of these deals --- is not established by these tests. We leave that mechanism open.

Overall, we read the result as supported but fragile. One formal test at the five-percent level supports H1a's claim that the run-up concentrates in cash deals --- but its significance rides on the imprecise negative stock estimate ($-0.0524$, standard error $0.0436$), and the cause's own cash-sidedness is weaker still (marginal, and only on the full panel). For that reason we keep our wording at \emph{concentrated} rather than \emph{specific}, we do not treat the difference between a significant and an insignificant coefficient as the test, and we assert no cause.

\section{Additional Analyses and Robustness}

\subsection{4.1}

The most immediate rival to our reading is analyst scrutiny. The idea is that the quarters just before a deal draw harder analyst questions about cash, and that it is the hard questions, not the deal itself, that produce the hedged answers we measure. This confound is genuine, not a straw man. Column~4 of Table~\ref{tab:empire_building_did} shows that the incidence of cash scrutiny --- whether a call draws heavy cash questioning at all --- does itself rise in the quarter before a cash acquisition: the coefficient on $\mathrm{HighCashScrutiny}$ is $0.0408^{**}$ (standard error $0.0194$). Scrutiny therefore moves together with both the firm's cash position and the event, and could in principle generate the run-up mechanically. The continuous, volume-based measure tells a weaker story --- in column~3 the $\mathrm{CashScrutiny}$ coefficient is $0.0854$ (standard error $0.0749$), not significant --- so the genuine confound is whether cash scrutiny happens at all (its incidence), not how much of a call it fills (its volume).

We rule it out in three steps. First, we confirm that the scrutiny measure is real --- that it actually moves with how much cash a firm holds. Table~\ref{tab:cash_scrutiny_validity} shows that it does. The share of analyst question-and-answer turns spent on cash and liquidity rises strongly with the firm's cash-to-assets ratio: $0.7530^{***}$ (standard error $0.0949$, one-tailed) without controls, and $0.8519^{***}$ (standard error $0.0967$) once the full control set is added. It also rises with a high-cash indicator, a simple flag for cash-rich firms, $0.1754^{***}$ (standard error $0.0237$). The measure clearly registers how cash-rich a firm is, so the nulls that follow are not the artifact of a dead variable that fails to pick anything up.

Second, scrutiny on its own does not move the residual --- our measure of how hedged a CEO's answers are, taken net of that CEO's persistent speaking style ($\mathrm{UncResCEO}$). In Table~\ref{tab:cash_scrutiny_channel}, the direct effect of scrutiny on that residual is zero to four decimal places (Panel~A, $-0.0000$, standard error $0.0013$, across 41,512 calls), and the same null holds when we run a logit on a dichotomized version of the residual --- a simple high-versus-low split (Panel~B, $-0.0003$, standard error $0.0074$) --- as well as on the raw uncertainty measure and on coarser scrutiny proxies. This test is not circular --- it is not merely the residual predicting itself --- because $\mathrm{CashScrutiny}$ counts the \emph{volume} of cash questions, whereas the uncertainty \emph{wording} of those questions ($\mathrm{UncQue}$) is already netted out of the residual when it is built, in the decomposition of Section~2.3.

Third, and most directly, we put scrutiny and the pre-announcement indicator into the same model, estimated on the cash-acquirer pre-announcement quarters. In Table~\ref{tab:reason_gating}, the pre-announcement indicator stays significant --- $0.0413^{***}$ (standard error $0.0177$) with scrutiny entered alongside it, and $0.0439^{**}$ (standard error $0.0189$) once the interaction is added --- while scrutiny's own effect is null ($0.0000$ and $0.0001$ across the two columns) and the scrutiny-by-pre-announcement interaction is insignificant ($-0.0056$, standard error $0.0111$). Put plainly: the reason for the deal raises uncertainty, the questioning does not, and the questioning does not amplify the reason either. The comparison is within firm --- each treated firm measured against its own ordinary quarters --- and correlational.

We keep the register honest about what this does and does not show. What we have is a failure to find an effect, paired with a dissociation inside a single model --- not a powered equivalence test that could formally rule an effect out. The interaction's confidence interval (roughly $[-0.027,\,+0.016]$) does not exclude effects that would still matter against a run-up near $0.04$; about 89\% of calls draw no cash scrutiny at all; and a contrast between a significant and an insignificant coefficient is not, by itself, a test. So our claim is the narrow one: analyst scrutiny does not account for \emph{this} run-up. We do not claim that scrutiny never matters anywhere. The reading is correlational and supportive, not definitive; it identifies no cause, and the mechanism remains open.

\subsection{4.2}

We now ask whether the speech-uncertainty signal moves the post-call information environment --- how large the information gap is between the firm and outside traders. We proxy that environment with the firm's bid-ask spread, the gap between the price to buy and the price to sell, measured over the 25 trading days that begin on the call date. This is a standard measure of how much outside traders fear trading against better-informed insiders, following BGT, whose post-call window we adopt, and we decompose the spread by the three speech components of Section~2.3. Two priors frame the test. DWZ found that the residual uncertainty is unrelated to stock-price or trading-volume responses and explains little of the market reaction, which loads instead on the chief executive's persistent clarity --- but they never examined a liquidity outcome. BGT, by contrast, found that call language relates to information asymmetry with \emph{opposite} signs across the call's two segments: complexity in the scripted presentation is positively associated with information asymmetry (an obfuscation reading), while complexity in the spontaneous response is negatively associated with it (an information reading).

First, the residual is inert. Across all twelve specifications in Table~\ref{tab:h14c_ceo2_decomp}, UncResCEO carries no association with the post-call spread. In the six contemporaneous columns the coefficients are imprecisely estimated and insignificant, ranging from $-0.0594$ to $-0.1021$ (42,625 calls); one quarter ahead they remain insignificant, ranging from $0.0545$ to $0.0984$ (43,049 calls). Because DWZ show this residual has no detected link to stock price or trading volume, and never test a liquidity outcome, this extends that finding to the bid-ask spread; the absence of an association here is not evidence of a precisely estimated zero, but it is what we find, and it sits comfortably with BGT's result that the spontaneous response segment does not widen information asymmetry.

Second, the scripted presentation is the segment outsiders appear to react to. UncPreCEO is positively associated with the contemporaneous spread, and significantly so in four of the six contemporaneous columns ($0.1664^{**}$, standard error $0.0946$, column~1; $0.3496^{***}$, standard error $0.1253$, column~2; $0.2945^{***}$, standard error $0.1229$, column~4; and $0.1644^{*}$, standard error $0.1184$, column~6; one-tailed). These four significant columns include all three firm-fixed-effect specifications (comparisons made within each firm). UncPreCEO is insignificant in the two industry-fixed-effect specifications with extended controls (columns~3 and~5, $0.0814$ and $0.0108$) and in all six one-quarter-ahead columns, and ClarityCEO is null throughout all twelve. The pattern is consistent with BGT's positive relation between scripted-presentation language and information asymmetry (the obfuscation component), and our own bid-ask result locates the contemporaneous spread with the scripted presentation rather than the residual. It does not contradict DWZ's finding that the residual is unrelated to stock price and volume, because DWZ never test the presentation and the bid-ask spread is an information-asymmetry outcome, not a price or volume reaction. We note that this association is contemporaneous only.

Taken together, the contemporaneous spread is associated with the scripted presentation but not with the residual; we do not test the between-segment difference directly. This pattern sits within DWZ's own scripted-versus-spontaneous framing and matches BGT's presentation-versus-response split. The payoff for our main analysis is one of interpretation rather than identification: because outsiders do not react to the residual, the pre-announcement run-up documented in Section~3 is unlikely to be an artifact of outsider reaction, and is more readily read as a private, insider-side signal. We offer this as a supportive reading, not as proof.

\subsection{4.3}

One feature of the Section~3.3 event study deserves a closer look. The post bin is dated to completion: a deal's quarters enter the post bin only once it has closed, and the quarters following a withdrawal are dropped from the event clock. Two of the three round-trip contrasts, peak-to-completion and gap-to-completion, are therefore estimated only over deals that go on to complete. Because completing and withdrawn deals need not be alike, part of the post-completion decline could in principle reflect \emph{which} deals complete rather than the resolution of the disclosure state itself. This is a correlational concern, and we do not claim to identify it away. The peak-to-gap contrast, however, is exempt. It is measured at the public announcement --- a point every announced deal reaches whatever its eventual outcome --- so it conditions on nothing that follows.

The check itself is straightforward. A withdrawal ends the disclosure state as surely as a completion does: once a pending deal is called off and disclosed, the executive is no longer answering around it. So we redefine the post bin to let the quarters at and after a withdrawal enter it alongside the completed quarters, and we re-estimate the Section~2.4 event study on this resolution-inclusive sample, reported in Table~\ref{tab:empire_drop_resolution}. Folding in the withdrawn deals asks a simple question: is the post-resolution decline particular to the deals that complete, or does it appear whenever the disclosure state resolves? We keep the framing modest. We do not claim that a weaker decline here would prove the main result an artifact; the power of this check is limited, as we discuss below. The reading is correlational and supportive rather than definitive.

Table~\ref{tab:empire_drop_resolution} shows that the round trip survives. On the resolution-inclusive sample the residual still falls from its pre-announcement peak to the resolved quarters by $0.0687$ ($p<.01$), against $0.0723$ ($p<.01$) in the main specification. The announcement-quarter leg, from peak to gap, is $0.0457$ ($p<.05$), almost exactly the $0.0455$ ($p<.05$) of the main result. And the cash leg declines from its pre-announcement level to resolution by $0.0204$ ($p<.01$). The sample here covers 28,191 firm-quarters over 1,321 firms. Treating a withdrawn deal as a resolution leaves the round trip in place.

We bound the check honestly. Few first deals are withdrawn and then observed for complete-case quarters inside the post window, so the resolution-inclusive sample adds only 89 firm-quarters and one firm to the 28,102 firm-quarters over 1,320 firms of the main analysis. With so little new data, the post-completion decline can barely move, so what this shows is that the handful of withdrawn deals behave like the completed ones --- not that a completion-driven reading is ruled out on its own. The firmer reassurance is structural: the peak-to-gap leg is measured at announcement and conditions on nothing that follows, so the resolution of the residual at announcement never depended on which deals complete. We therefore read the resolution-inclusive estimates as consistent with the main specification, not as an independent test.

\subsection{4.4}

The cash-ratio regressions carry the cash ratio's own one-quarter lag alongside firm fixed effects --- a partial-adjustment form, well suited to a sticky stock like cash that moves toward its target only gradually. A dynamic fixed-effects panel of this kind is subject to the Nickell bias, a small-sample bias that the within-firm transformation induces. That bias is not confined to the lag coefficient; it also reaches regressors correlated with lagged cash, so simply declining to interpret the lag does not immunize the event-time bins. The clean guard is therefore to remove the dynamic term altogether, rather than to reason about the size of the bias. The residual takes no lag and is untouched by this, so it serves as our anchor. This is a specification concern, not a claim of causal identification.

So we re-estimate the cash column with static fixed effects, dropping the lagged dependent variable (Table~\ref{tab:empire_drop_staticfe}). Removing the lag changes what the coefficients describe. With the lag, each coefficient is a \emph{change} in the cash ratio around its own recent level --- a partial adjustment. Without it, the coefficients describe the \emph{level} of the cash ratio relative to the baseline. The two specifications are on different scales, so the comparison that carries over is the timing pattern --- the sign and significance of the bins across event time --- not the size of any single coefficient. This is a robustness check and is supportive rather than definitive.

Table~\ref{tab:empire_drop_staticfe} shows that the timing pattern carries over. Without the lag, the cash ratio still does not fall when the deal is announced: the pre-announcement-to-gap change is $0.0012$ below zero and insignificant, much like the $0.0006$ of the with-lag specification. It falls only once the cash leaves the firm at completion, with a gap-to-completion drop of $0.0318$ ($p<.01$) and a completion coefficient of $0.0189$ below zero ($p<.01$). The pre-announcement-to-completion decline is $0.0305$ ($p<.01$), larger than the $0.0216$ ($p<.01$) of the partial-adjustment specification, exactly as expected when the coefficients describe a level rather than a change. The residual column is unchanged, since it never carried the lag, and is identical to Table~\ref{tab:empire_drop_matched}.

How should we read the two side by side? The lag belongs in the main specification, because cash is sticky and partial adjustment is the right description of how it moves; the static specification is not a better model but a check that removes the dynamic term entirely. And because the static specification carries no lagged dependent variable, it carries no dynamic-panel bias --- so the cash timing conclusion (no fall at announcement, a fall at completion) holds whether or not that term is present. The timing of the cash path, on which the differential-timing result rests, does not ride on the dynamic specification. The reading is correlational and supportive rather than definitive.

\subsection{4.5}

Our last robustness check drops the first-deal restriction and stacks every deal a firm makes. The pre-announcement run-up survives. On the all-deals-stacked panel of Table~\ref{tab:empire_building_did}, the cash-arm residual-uncertainty coefficient is $0.0391^{***}$ (standard error $0.0140$) and the cash-ratio coefficient is $0.0033^{**}$ (standard error $0.0015$), while the stock-arm residual coefficient is $-0.0348$ (standard error $0.0272$), a flat null. It also reappears in a binary forward form. Among all firm-quarters, higher CEO Q\&A residual uncertainty is associated with a deal being announced the next quarter, of any payment type: a linear-probability coefficient of $0.0086^{***}$ (standard error $0.0026$, $p=.0011$) and a logit coefficient of $0.3233^{***}$ (standard error $0.0967$, $p=.0008$) in this forward specification (Logit~A), estimated on 40,004 firm-quarters from 1,422 firms with a deal rate of 2.84\% (pseudo-$R^2$ $0.026$). It survives firm and year-quarter fixed effects, with a linear-probability coefficient of $0.0078^{***}$ (standard error $0.00275$; 39,557 firm-quarters, 1,422 firms; within-$R^2$ $0.003$).

The differential-timing round trip survives all-deals as well. On the matched universe, Table~\ref{tab:empire_drop_matched} shows residual uncertainty spiking at PRE1 ($0.0352^{**}$, standard error $0.0148$) and then unwinding: the PRE1-to-GAP drop is $0.0363^{**}$ (standard error $0.0156$) and the PRE1-to-POST drop is $0.0544^{***}$ (standard error $0.0139$). Split by payment type, Table~\ref{tab:empire_drop_placebo} shows the round trip holding in the cash arm (PRE1 $0.0401^{***}$, standard error $0.0145$; PRE1-to-POST drop $0.0543^{***}$, standard error $0.0137$) and not in the stock arm (PRE1 $-0.0272$, standard error $0.0266$, not significant), where the stock arm's PRE1-to-GAP change is $-0.0585^{*}$ (standard error $0.0355$) --- marginally significant and negative, which we do not over-read given the noisy stock null.

The formal cash-concentration result also survives, and in fact strengthens. On the all-deals-stacked panel of Table~\ref{tab:empire_cashspec}, the pooled cash-minus-stock Wald difference is $0.1056^{**}$ (standard error $0.0423$; $z\approx2.50$, $p\approx.013$) --- the cash arm at $0.0447^{**}$ (standard error $0.0180$) against the stock arm's insignificant $-0.0609$ (standard error $0.0385$) --- larger than the first-deal difference of $0.0983^{**}$ ($p=.039$) in the same table. The proposed cause remains the weak leg: the matched cash-ratio cash-minus-stock difference is $0.0071$ (standard error $0.0062$), not significant. So the run-up still concentrates in cash deals --- concentration, not strict specificity, and with no cause asserted --- while the build-up of cash itself is not differentially cash-sided here. In binary form, among deals at $e=-1$, higher residual uncertainty is associated with the deal being a cash deal (coded one) rather than a stock deal (coded zero): a linear-probability coefficient of $0.0613^{**}$ (standard error $0.0283$, $p=.030$) and a logit coefficient of $0.7478^{**}$ (standard error $0.3399$, $p=.028$) in this specification (Logit~B), on 1,105 deal observations from 563 firms (982 cash, 123 stock; cash base rate 88.9\%; pseudo-$R^2$ $0.078$). Under firm and year-quarter fixed effects the coefficient keeps its sign but is no longer significant: a linear-probability coefficient of $0.0644$ (standard error $0.05076$, not significant; $p=.205$; 1,063 deal observations, 563 firms; within-$R^2$ $0.059$). The design stays correlational and identifies no causal effect, and we read it as concentration in cash deals, not strict specificity.

\section{5}

In this study we examine whether the language a chief executive uses in the unscripted question-and-answer session of an earnings call tracks the disclosure state of an acquisition the firm is still withholding, using a residual measure of chief-executive question-and-answer uncertainty that nets out each executive's persistent speaking style.

We find that residual uncertainty is elevated in the quarter before a cash acquisition, with no comparable rise before stock acquisitions; that it becomes indistinguishable from zero once the deal is announced, even before it closes, while the cash that funds it persists to completion; and that the elevation concentrates in cash acquisitions rather than stock on a formal pooled test.

The pattern is not accounted for by analysts devoting more of the call to cash questions; the pre-announcement elevation survives a direct measure of analyst cash scrutiny and its interaction with the pre-announcement window. A further analysis finds the residual unrelated to the firm's post-call bid-ask spread, while the scripted presentation is contemporaneously positively associated with it, consistent with the scripted and unscripted segments relating to the outside information environment differently. The differential-timing result is consistent with treating a withdrawal as a resolution of the deal, and it holds when the cash specification is estimated without its dynamic term; and the run-up and cash-concentration hold when every deal a firm makes is stacked.

Taken together, these patterns suggest that the unscripted language of earnings calls carries a readable, anticipatory trace of a deal's passage from private to public, a signature researchers can in principle observe before the announcement. For the academic literature, the pattern connects two strands usually studied apart: the theory of when an informed manager withholds private information, and the evidence that markets anticipate corporate transactions before they are announced. The within-firm language signature we characterize sits in the window between them. We draw no prescription for investors, managers, or regulators from a correlational pattern that identifies no causal channel; the contribution is to characterize a regularity, not to recommend acting on it.

These findings carry clear limitations. The evidence is correlational and within-firm, and the design supports no causal identification and establishes no mechanism. The concentration in cash deals is motivated in Section~2 but not identified, and the war-chest channel behind cash accumulation in particular remains unestablished. Separately, the source of the uncertainty, a compliance-constrained inability to speak versus a strategically chosen reticence, remains observationally equivalent.

Two further boundaries qualify the evidence. The residual is estimated only for chief executives with enough calls to fix a speaking style, which skews the language sample toward larger, more heavily covered firms, and the scrutiny rule-out is a failure to find rather than a powered test of equivalence, since most calls draw no cash scrutiny at all. The measure itself bounds the inference further. Uncertainty is captured by applying a finance-specific word list to the transcripts, a count of words that abstracts from their context, and the residual is one operationalization of call-specific uncertainty rather than a direct reading of what a chief executive knows. The stock-financed comparison shares the disclosure bind but is an imperfect counterfactual, and the sample, United States public firms from 2002 to 2018, approximately the S\&P~1500, may not extend to other periods, other markets, or the smaller firms outside it.

Future work could ask whether the same anticipatory language signature appears around other classes of withheld material events, and could pursue designs able to separate compliance-constrained from strategic silence, which our setting cannot distinguish. Several extensions follow. The same residual-uncertainty reading could be applied to other corporate transactions, such as divestitures, joint ventures, and strategic alliances, and to settings outside the United States; richer measures of spoken uncertainty could replace the word-list count; and identifying the mechanism behind the cash concentration, which Section~2 motivates but does not establish, including the cash-accumulation channel this study leaves open, would require a design the present setting does not provide.

\clearpage
\begin{table}[htbp]\centering\caption{Cash Scrutiny Channel}\label{tab:cash_scrutiny_channel}\textit{[table content]}\end{table}

\begin{table}[htbp]\centering\caption{Cash Scrutiny Validity}\label{tab:cash_scrutiny_validity}\textit{[table content]}\end{table}

\begin{table}[htbp]\centering\caption{Empire Building Did}\label{tab:empire_building_did}\textit{[table content]}\end{table}

\begin{table}[htbp]\centering\caption{Empire Cashspec}\label{tab:empire_cashspec}\textit{[table content]}\end{table}

\begin{table}[htbp]\centering\caption{Empire Drop Matched}\label{tab:empire_drop_matched}\textit{[table content]}\end{table}

\begin{table}[htbp]\centering\caption{Empire Drop Placebo}\label{tab:empire_drop_placebo}\textit{[table content]}\end{table}

\begin{table}[htbp]\centering\caption{Empire Drop Resolution}\label{tab:empire_drop_resolution}\textit{[table content]}\end{table}

\begin{table}[htbp]\centering\caption{Empire Drop Staticfe}\label{tab:empire_drop_staticfe}\textit{[table content]}\end{table}

\begin{table}[htbp]\centering\caption{H14C Ceo2 Decomp}\label{tab:h14c_ceo2_decomp}\textit{[table content]}\end{table}

\begin{table}[htbp]\centering\caption{Reason Gating}\label{tab:reason_gating}\textit{[table content]}\end{table}

\begin{table}[htbp]\centering\caption{Summary Stats}\label{tab:summary_stats}\textit{[table content]}\end{table}

\begin{thebibliography}{99}

\bibitem[Basic Inc.\ v.\ Levinson(1988)]{basic_v_levinson}
Basic Inc.\ v.\ Levinson, 485 U.S. 224 (1988).

\bibitem[Bates et~al.(2009)]{bates2009}
Bates, T.~W., K.~M. Kahle, and R.~M. Stulz. 2009. Why do U.S.\ firms hold so much more cash than they used to? \emph{The Journal of Finance} 64: 1985--2021. % VERIFY pages

\bibitem[Bertrand and Schoar(2003)]{bertrand_schoar2003}
Bertrand, M., and A.~Schoar. 2003. Managing with style: The effect of managers on firm policies. \emph{The Quarterly Journal of Economics} 118: 1169--1208.

\bibitem[Dzielinski et~al.(2021)]{dwz}
Dzielinski, M., A.~F. Wagner, and R.~J. Zeckhauser. 2021. Straight talkers and vague talkers: The effects of managerial style in earnings conference calls. Working paper. % alias of dwz2021 (props use both keys -> normalize to one in the harness)

\bibitem[Dzielinski et~al.(2021)]{dwz2021}
Dzielinski, M., A.~F. Wagner, and R.~J. Zeckhauser. 2021. Straight talkers and vague talkers: The effects of managerial style in earnings conference calls. Working paper. % VERIFY year/venue (working paper)

\bibitem[Dye(1985)]{dye1985}
Dye, R. 1985. Disclosure of nonproprietary information. \emph{Journal of Accounting Research} 23: 123--145.

\bibitem[Harford(1999)]{harford1999}
Harford, J. 1999. Corporate cash reserves and acquisitions. \emph{The Journal of Finance} 54: 1969--1997.

\bibitem[Hollander et~al.(2010)]{hollander2010}
Hollander, S., M.~Pronk, and E.~Roelofsen. 2010. Does silence speak? An empirical analysis of disclosure choices during conference calls. \emph{Journal of Accounting Research} 48: 531--563.

\bibitem[Keown and Pinkerton(1981)]{keown1981}
Keown and Pinkerton. 1981. Merger announcements and insider trading activity: An empirical investigation. \emph{The Journal of Finance} 36: 855--869.

\bibitem[Loughran and McDonald(2011)]{lm2011}
Loughran, T., and B.~McDonald. 2011. When is a liability not a liability? Textual analysis, dictionaries, and 10-Ks. \emph{The Journal of Finance} 66: 35--65.

\bibitem[Louis(2004)]{louis2004}
Louis, H. 2004. Earnings management and the market performance of acquiring firms. \emph{Journal of Financial Economics} 74: 121--148.

\bibitem[Matsumoto et~al.(2011)]{matsumoto2011}
Matsumoto, D., M.~Pronk, and E.~Roelofsen. 2011. What makes conference calls useful? The information content of managers' presentations and analysts' discussion sessions. \emph{The Accounting Review} 86: 1383--1414.

\bibitem[Opler et~al.(1999)]{opler1999}
Opler, T., L.~Pinkowitz, R.~Stulz, and R.~Williamson. 1999. The determinants and implications of corporate cash holdings. \emph{Journal of Financial Economics} 52: 3--46. % VERIFY pages

\bibitem[Pagan(1984)]{pagan1984}
Pagan, A. 1984. Econometric issues in the analysis of regressions with generated regressors. \emph{International Economic Review} 25: 221--247. % VERIFY pages

\bibitem[Ragozzino(2024)]{ragozzino2024}
Ragozzino, R. 2024. [TITLE TK]. \emph{Long Range Planning} [vol./pages TK]. % SOURCE: notebook PII S0024-6301(23)00100-0 -- SINA VERIFY authors/title/pages

\bibitem[U.S.\ SEC(Rule 10b-5)]{rule_10b5}
U.S.\ Securities and Exchange Commission. Rule 10b-5: Employment of manipulative and deceptive devices. 17 C.F.R.\ \S\,240.10b-5.

\bibitem[Shleifer and Vishny(2003)]{shleifer_vishny2003}
Shleifer, A., and R.~Vishny. 2003. Stock market driven acquisitions. \emph{Journal of Financial Economics} 70: 295--311.

\bibitem[Thewissen et~al.(2024)]{thewissen2024}
Thewissen, J., et~al. 2024. [TITLE TK -- tone working paper]. SSRN working paper 4900453. % SINA VERIFY authors/title

\bibitem[Verrecchia(1983)]{verrecchia1983}
Verrecchia, R. 1983. Discretionary disclosure. \emph{Journal of Accounting and Economics} 5: 179--194.

\end{thebibliography}
\end{document}
